% formulalist.tex: Must-know relations for the Planetary Systems exam
% Planetary Systems, Kapteyn Institute, University of Groningen

\documentclass[11pt]{article}
\usepackage{../../worksheets/common/worksheetIPS}

% This document is not a worksheet: replace the running head accordingly.
\fancyhead[L]{\footnotesize\color{ipsPrimary}Planetary Systems \textbullet\ Exam formula list}
\fancyhead[R]{\footnotesize\color{ipsPrimary}Must-know relations}

% Numbered relation entry: title line, displayed equation, short explanation.
\newcounter{ipsrelation}
\newcommand{\relation}[1]{%
  \par\addvspace{11pt plus 3pt}%
  \refstepcounter{ipsrelation}%
  \noindent{\large\bfseries\color{ipsPrimary}\arabic{ipsrelation}.\enspace #1\par}%
  \vskip1pt
  \noindent{\color{ipsAccent}\rule{\linewidth}{0.5pt}}%
  \par\addvspace{3pt}%
}

% Lecture links point at the published course site.
\newcommand{\lecref}[2]{\href{https://ips.formingworlds.space/#1}{#2}}

% Total derivative, matching the lecture-notes macro of the same name.
\providecommand{\dv}[2]{\frac{\mathrm{d} #1}{\mathrm{d} #2}}

\begin{document}
\thispagestyle{fancy}

\begingroup
\noindent{\color{ipsAccent}\rule{\linewidth}{1.2pt}}\\[8pt]
{\LARGE\bfseries\color{ipsPrimary} Exam formula list: must-know relations\par}
\vskip7pt
{\normalsize\color{ipsAccent} Planetary Systems final exam \textbullet\ closed book\par}
\vskip3pt
{\footnotesize Planetary Systems (WBAS002-05) \textbullet\ Kapteyn Astronomical Institute, University of Groningen\par}
\vskip7pt
\noindent{\color{ipsAccent}\rule{\linewidth}{0.6pt}}
\vskip10pt
\endgroup

The final exam is closed book: a calculator and a pen are all that is allowed.
The twelve relations on this list are the ones you are expected to know from memory and apply without prompting.
Every other formula a question needs, for example the vis-viva equation, the Roche limit, or the Jeans escape flux, is printed in the question itself.
Each entry below states where the relation is introduced in the \href{https://ips.formingworlds.space}{lecture notes}; the same list with the same links appears on the \href{https://ips.formingworlds.space/mock_exams.html}{mock exams page}.

\relation{Newton's law of gravitation}
\begin{equation*}
F = \frac{G M m}{r^2}, \qquad U = -\frac{G M m}{r}
\end{equation*}
The attractive force between masses $M$ and $m$ at separation $r$, and the corresponding gravitational potential energy.
This pair underlies all orbital dynamics and the energy release of accretion and differentiation.
The force law drives the orbit derivations of \lecref{01_introduction/introduction.html}{Lecture 1}; the energy form is introduced in \lecref{02_formation_orbits/formation_orbits.html}{Lecture 2} and drives the heating budgets of \lecref{04_differentiation_magnetospheres/differentiation_magnetospheres.html}{Lecture 4}.

\relation{Circular orbit speed}
\begin{equation*}
v_{\mathrm{c}} = \sqrt{\frac{G M}{r}}
\end{equation*}
The speed of a body on a circular orbit of radius $r$ around a central mass $M$, from the balance of gravitational and centripetal acceleration.
Derived in \lecref{01_introduction/introduction.html}{Lecture 1}; generalised by the vis-viva equation of \lecref{02_formation_orbits/formation_orbits.html}{Lecture 2}.

\relation{Escape velocity}
\begin{equation*}
v_{\mathrm{esc}} = \sqrt{\frac{2 G M}{r}}
\end{equation*}
The minimum speed needed to escape to infinity from distance $r$, from setting kinetic plus potential energy to zero.
It controls impact energetics and which gases a planet can retain.
Introduced in \lecref{02_formation_orbits/formation_orbits.html}{Lecture 2}; central to atmospheric escape in \lecref{05_atmospheres_1/atmospheres_1.html}{Lecture 5}.

\relation{Kepler's third law}
\begin{equation*}
P^2 = \frac{4 \pi^2 a^3}{G M}
\end{equation*}
The orbital period $P$ of an orbit with semi-major axis $a$ around a central mass $M$.
It is the primary tool for weighing central bodies and for converting exoplanet periods into orbital distances.
Derived in Newtonian form in \lecref{01_introduction/introduction.html}{Lecture 1}; restated in \lecref{02_formation_orbits/formation_orbits.html}{Lecture 2} and applied to exoplanets in \lecref{13_exoplanets/exoplanets.html}{Lecture 13}.

\relation{Mean density}
\begin{equation*}
M = \frac{4}{3} \pi R^3 \bar{\rho}
\end{equation*}
The relation between mass, radius, and bulk density $\bar{\rho}$.
Bulk density is the first-order constraint on whether a body is made of rock, ice, or gas.
Used throughout \lecref{08_interiors/interiors.html}{Lecture 8} and \lecref{13_exoplanets/exoplanets.html}{Lecture 13}.

\relation{Hydrostatic equilibrium}
\begin{equation*}
\dv{P}{r} = -\rho g
\end{equation*}
The pressure gradient that balances gravity in any static fluid layer.
It structures both atmospheres and interiors and is the starting point of the scale-height and central-pressure estimates.
Introduced in \lecref{05_atmospheres_1/atmospheres_1.html}{Lecture 5} and \lecref{08_interiors/interiors.html}{Lecture 8}.

\relation{Ideal gas law}
\begin{equation*}
P = n \kB T = \frac{\rho \kB T}{\mu\, m_u}
\end{equation*}
The equation of state of a dilute gas, with number density $n$, mean molecular weight $\mu$, and atomic mass unit $m_u$.
It closes the hydrostatic equation for atmospheres.
Introduced in \lecref{05_atmospheres_1/atmospheres_1.html}{Lecture 5}.

\relation{Isothermal scale height}
\begin{equation*}
H = \frac{\kB T}{\mu\, m_u\, g}
\end{equation*}
The e-folding height of pressure in an isothermal atmosphere, from hydrostatic equilibrium plus the ideal gas law.
Hotter or lighter atmospheres are more extended; stronger gravity compresses them.
Derived in \lecref{05_atmospheres_1/atmospheres_1.html}{Lecture 5}.

\relation{Stefan--Boltzmann law}
\begin{equation*}
F = \sigma T^4
\end{equation*}
The radiative flux from an ideal (blackbody) surface at temperature $T$, with $\sigma$ the Stefan--Boltzmann constant.
\lecref{03_heat_energy/heat_energy.html}{Lecture 3} states the general form $F = \epsilon \sigma T^4$, with emissivity $\epsilon \leq 1$ close to 1 for most planetary surfaces.
It governs how planets shed heat to space.

\relation{Equilibrium temperature (energy balance)}
\begin{equation*}
\frac{S}{4} (1 - A) = \sigma T_{\mathrm{eq}}^4
\qquad \Longrightarrow \qquad
T_{\mathrm{eq}} = \left[ \frac{S (1 - A)}{4 \sigma} \right]^{1/4}
\end{equation*}
Absorbed stellar flux (stellar constant $S$, Bond albedo $A$, the factor 4 from the ratio of a sphere's surface to its cross-section) equals emitted thermal flux.
This balance sets the baseline temperature of every planet before greenhouse warming.
Derived in \lecref{05_atmospheres_1/atmospheres_1.html}{Lecture 5}; extended to the greenhouse model in \lecref{09_earth_venus/earth_venus.html}{Lecture 9} and applied to exoplanets in \lecref{13_exoplanets/exoplanets.html}{Lecture 13}.

\relation{Thermal energy per particle}
\begin{equation*}
E_{\mathrm{th}} = \frac{3}{2} \kB T
\end{equation*}
The mean kinetic energy of a gas particle at temperature $T$.
Comparing thermal and gravitational energy per molecule decides between retention and escape of an atmosphere.
The escape physics of \lecref{05_atmospheres_1/atmospheres_1.html}{Lecture 5} and \lecref{10_mercury_mars/mercury_mars.html}{Lecture 10} builds on this comparison; the Jeans parameter there uses $\kB T$ as the thermal scale.

\relation{Radioactive decay}
\begin{equation*}
N(t) = N_0\, e^{-\lambda t}, \qquad t_{1/2} = \frac{\ln 2}{\lambda}
\end{equation*}
The exponential decay of a parent isotope with decay constant $\lambda$ and half-life $t_{1/2}$.
It powers radiogenic heating of interiors and underlies all radiometric dating.
The half-lives behind radiogenic heating appear in \lecref{03_heat_energy/heat_energy.html}{Lecture 3}; the decay law itself is introduced with radiometric dating in \lecref{12_small_bodies/small_bodies.html}{Lecture 12}.

\end{document}
